% ---------------------------------------------------------------------
\begin{tcolorbox}[enhanced,breakable,title=常见误解,fonttitle=\bfseries,
  colback=gray!3!white,colframe=gray!55!black,boxrule=0.6pt,arc=2pt,
  left=5pt,right=5pt,top=3pt,bottom=3pt,before skip=6pt,after skip=6pt]
\textbf{“$B\to\pi\pi$ 的 QCD 因子化已对所有阶证明。”}不对：两个轻介子末态的领头幂次公式只在有限阶确立，分歧未决（\cref{prop:heavy-qcdf-lp}，[P]）；[A] 只属于 $B\to D\pi$ 与重到轻形状因子。\textbf{“BBNS 公式只在 $\order{\as}$ 成立。”}也不对：它是对 $\as$ 所有阶的领头幂次陈述，高阶检验与 SCET 推导都支持它。\par
\textbf{“$X_A$、$X_H$ 的参数化是因子化定理的一部分。”}不对：它们参数化发散的卷积（\cref{prop:heavy-qcdf-endpoint}，[V]），是模型。\par
\textbf{“NRQCD 衰变因子化与 OPE 同样严格。”}不对：衰变率是固定物理点上的 Minkowski 量，非微扰地需要局域对偶（\cref{conj:ee-duality}(b)）。\par
\textbf{“NRQCD 产生的因子化已经证明；NQS 的 Wilson 线在 NNLO 不改变结果。”}都不对：产生因子化是猜想（\cref{conj:heavy-nrqcd-production}），规范完备的 Wilson 线在 NNLO 是必需的。
\end{tcolorbox}

\paragraph{本章小结。}
以重夸克质量为硬标度的因子化中，只有电磁衰变的匹配系数在有限阶是严格的（[R]）；$B\to D\pi$、重到轻形状因子、单举 $B$ 衰变、NRQCD 衰变与大 $\pT$ 偶素的碎裂因子化有全阶论证（[A]）；两个轻介子末态只在有限阶确立（[P]），其幂次修正的标准因子化失效（[V]）；NRQCD 产生是猜想（[C]）。用于固定运动学点上的强子衰变率，这些结果都还需要局域对偶（[C]）。
